\section{Open Data, Software Codebases, and Experimental Specifications}
\label{sec:data_code_spec}

\subsection{Summary of Resources and Infrastructure}

\begin{table*}[!h]
\centering
\footnotesize
\renewcommand{\arraystretch}{0.88}
\setlength{\tabcolsep}{4pt}
\caption{Summary of Open Datasets, Audited Software Codebases, and Hardware Infrastructure.}
\label{tab:key_resources}
\begin{tabularx}{\textwidth}{p{4.0cm}p{5.5cm}>{\raggedright\arraybackslash}X}
\toprule
\textbf{Category} & \textbf{Resource / Data Source} & \textbf{Identifier / Access Link} \\
\midrule
\multicolumn{3}{l}{\textbf{Commercial Wind Farm Telemetry and Market Clearing Records}} \\
\midrule
Kelmarsh Wind Farm 10-min SCADA & Cubico Sustainable Investments / Zenodo & DOI: \href{https://doi.org/10.5281/zenodo.5841834}{10.5281/zenodo.5841834} \\
Penmanshiel Wind Farm 10-min SCADA & Cubico Sustainable Investments / Zenodo & DOI: \href{https://doi.org/10.5281/zenodo.5946808}{10.5281/zenodo.5946808} \\
Hill of Towie Wind Farm 10-min SCADA & TRIG / RES / Zenodo (Open Data) & DOI: \href{https://doi.org/10.5281/zenodo.20204946}{10.5281/zenodo.20204946}; Elexon BM Unit: \texttt{T\_HILW-1} \\
UK Elexon BMRS Market Clearing Prices & Elexon Balancing and Settlement Code & Open Government Licence v3.0 \\
Systematic Scoping Audit Corpus ($N=46$) & This Paper (Verbatim Extraction Ledger) & Data Repository: \texttt{corpus48\_consensus} \\
\midrule
\multicolumn{3}{l}{\textbf{Open-Source Software Implementations and Audit Scripts}} \\
\midrule
Python Programming Language (v3.10) & Python Software Foundation & RRID: \href{https://scicrunch.org/resolver/SCR_008394}{SCR\_008394} \\
PyTorch Deep Learning Library (v2.0+) & PyTorch Foundation & RRID: \href{https://scicrunch.org/resolver/SCR_018536}{SCR\_018536} \\
Statsmodels (Circular Block Bootstrap) & Statsmodels Developers & RRID: \href{https://scicrunch.org/resolver/SCR_016074}{SCR\_016074} \\
Scikit-learn (Machine Learning in Python) & Pedregosa et al., 2011 & RRID: \href{https://scicrunch.org/resolver/SCR_002577}{SCR\_002577} \\
DLinear Baseline Codebase & Zeng et al., AAAI 2023 & GitHub: \texttt{cure-lab/LTSF-Linear} \\
PatchTST Benchmark Codebase & Nie et al., ICLR 2023 & GitHub: \texttt{yuqinie98/PatchTST} \\
Audited PCT Codebase & Wang et al., \textit{Sci. Rep.} 2026 & GitHub: \texttt{wind-power-forecasting/PCT} \\
Audited PINN-UQ Codebase & Xu et al., \textit{Eng. App. AI} 2025 & GitHub: \texttt{xuchen/PINN-UQ} \\
Full Empirical Audit and Verification Suite & This Paper & GitHub: \href{https://github.com/anonymous-audit/pcwpf-audit}{anonymous-audit/pcwpf-audit} \\
\midrule
\multicolumn{3}{l}{\textbf{High-Performance Computing Hardware}} \\
\midrule
NVIDIA RTX 4090 Dedicated GPUs ($10\times$) & NVIDIA Corporation & Compute Capability 8.9 (24GB VRAM each) \\
\bottomrule
\end{tabularx}
\end{table*}

\subsection{Data and Code Availability}
Raw operational SCADA telemetry from the Kelmarsh (6 turbines, 12.3~MW, Senvion MM92) and Penmanshiel (14 turbines, 28.7~MW, Senvion MM82) wind farms is openly available on Zenodo under CC-BY 4.0 licenses from Cubico Sustainable Investments (DOIs: \href{https://doi.org/10.5281/zenodo.5841834}{10.5281/zenodo.5841834} and \href{https://doi.org/10.5281/zenodo.5946808}{10.5281/zenodo.5946808}). Operational SCADA telemetry from the Hill of Towie wind farm (21 turbines, 48.3~MW, Siemens SWT-2.3-82) is openly available on Zenodo under CC-BY 4.0 license released by RES on behalf of The Renewables Infrastructure Group (TRIG) (DOI: \href{https://doi.org/10.5281/zenodo.20204946}{10.5281/zenodo.20204946}). Elexon System Buy Price (SBP) and System Sell Price (SSP) records are freely accessible through the Elexon BMRS portal under the UK Open Government Licence v3.0. The complete, reproducible Python verification suite---including computational graph static parsers, 2,000-replication circular block bootstrap scripts, Elexon BSC physical settlement calculators, and figure reproduction scripts---is openly archived at Zenodo and GitHub (\href{https://github.com/anonymous-audit/pcwpf-audit}{\texttt{anonymous-audit/pcwpf-audit}}). The full 46-study scoping extraction database with verbatim textual quotes, equation numbers, and multi-coder Fleiss' $\kappa$ consensus scores is provided as open supplementary data.

Further technical information, trained model weights, or raw inter-rater reliability extraction matrices will be fulfilled by the corresponding author, Junyu Li (\href{mailto:ljylikezmn999@gmail.com}{\texttt{ljylikezmn999@gmail.com}}).

\subsection{Wind Farm Telemetry and Data Sanitation Protocols}
This computational study utilized field telemetry records from three commercial utility-scale wind farms located across the United Kingdom. The \textbf{Kelmarsh Wind Farm} (Northamptonshire, England) represents flat open inland topography and comprises six Senvion MM92 turbines (rated 2.05~MW each, hub height 78~m, rotor diameter 92~m, farm capacity 12.3~MW), with training observations spanning $N_{\text{train}} = 5,098$ clean 10-minute intervals. The \textbf{Hill of Towie Wind Farm} (Drummuir, Scotland) represents complex upland terrain featuring micro-topographical wake interactions, comprising 21 Siemens SWT-2.3-82 turbines (rated 2.3~MW each, hub height 60~m, rotor diameter 82~m, farm capacity 48.3~MW), evaluated as a full-year independent spatiotemporal extrapolation benchmark with 7,599 forecast origins (45,594 forward settlement half-hours in 2017). The \textbf{Penmanshiel Wind Farm} (Scottish Borders) represents undulating upland terrain comprising 14 Senvion MM82 turbines (rated 2.05~MW each, farm capacity 28.7~MW), configured with a restricted training partition ($N_{\text{train}} = 1,583$ intervals) specifically to evaluate over-fitting and prior inductive bias under sample-restricted regimes (acknowledging that site topography and turbine models co-vary with sample size across wind farms). Standard data sanitation protocols removed periods of turbine curtailment, anemometer icing, and communication maintenance dropouts ($P \le 0$ when $v > v_{\text{cut-in}}$ or active power frozen for $>6$ consecutive steps).

\subsection{Method Details and Statistical Estimation}

\paragraph*{Systematic Scoping Review Protocol (PRISMA-ScR)}
Literature identification followed a structured protocol specified prior to full-text extraction, executing four Boolean title queries via the OpenAlex API ($Q_1$ through $Q_4$), complemented by manual keyword-directed discovery across Google Scholar, arXiv, and preprint servers for papers published between January 2023 and March 2026 (with rolling monitoring through July 2026), yielding 89 raw candidate records. Two-phase screening (title/abstract followed by full-text extraction) identified our core scoping cohort of 46 contemporary studies (38 peer-reviewed journal articles and 8 preprints, 2023--2026) subjected to exhaustive five-dimensional ($\mathcal{D}_1$--$\mathcal{D}_5$) extraction and inter-rater reliability coding. A targeted sensitivity search expanding retrieval to title, abstract, and author keywords across Scopus and IEEE Xplore (encompassing broader descriptors such as \texttt{"knowledge-guided"}, \texttt{"hybrid"}, and \texttt{"domain-informed"}) retrieved 14 additional candidate papers, with full-text inspection confirming zero instances of documented operational failure boundaries, specifically stress-testing the 0.0\% negative trade-off finding. Code availability and ablation metrics (95.7\% open-code deficit and 50.0\% confounded module ablations) are strictly bounded to the audited 46-study cohort with Wilson confidence intervals. Inter-rater reliability was benchmarked across a representative 20-study pilot by three independent raters (Coders A, B, and C) achieving Fleiss' $\kappa = 0.646\text{--}1.000$ across provenance, stage, and telemetry dimensions.

\paragraph*{Static Code and Computational Graph Audit}
Software codebases were audited via static AST parsing and line-by-line algebraic comparison against published manuscripts. For PCT (\textit{Sci. Rep.} 2026), the loss definition in \texttt{tst/loss.py} was decompiled into symbolic PyTorch operations, establishing that \texttt{F.kl\_div(p.log(), q)} computes $\sum q_i \log(q_i / p_i) = \text{KL}(q \parallel p)$ instead of the published $\text{KL}(p \parallel q)$. Algebraic tensor expansion of the empirical density kernel $K_p^\top K_w$ verified strict equivalence to a 2D product kernel density estimator.

\paragraph*{10-GPU Prospective Controlled Ablation Testbed and Information Budget Matrix}
Ablation experiments were executed across dedicated GPU infrastructure: Kelmarsh was evaluated across 4 random initialization seeds (42, 43, 44, 45); Hill of Towie and Penmanshiel across 3 random seeds each (42, 43, 44) for benchmark ablations, and across 5 random seeds (42, 123, 456, 789, 1024) for learning curve and information budget experiments. All models were implemented in PyTorch under CUDA. Neural backbones were trained using Adam with learning rate $10^{-3}$ and batch size 128 in accordance with the frozen benchmark protocol. Regularization weight $\lambda$ was evaluated across a predefined sensitivity and stress-testing grid ($\lambda \in \{0.01, 0.1, 1.0\}$), representing weak regularization ($\lambda=0.01$), moderate enforcement ($\lambda=0.1$), and strong stress testing ($\lambda=1.0$). Hyperparameters, model selection, and quantile calibration offsets were determined exclusively on the 2016 validation and calibration partitions; the 2017 test year was accessed only after model configurations and hyperparameter regimes were frozen, preventing test-set selection bias. For pseudo-scrambled controls, the empirical wind-to-power mapping was permuted along the wind axis, preserving identical loss formulations, network architectures, and regularization weights $\lambda$. To disentangle domain physical guidance from historical sample expansion, a controlled information-budget experiment at $N_{\text{train}}=500$ evaluated models against both budget-matched priors (fitted strictly on the identical 500 samples) and full-history priors (fitted on 7--10 months). The primary forecasting models (DLinear) take exclusively historical power sequences ($x[:, :, 0]$) as inputs at inference, requiring zero runtime meteorological telemetry. The historical wind speed proxy lag scan ($k \in \{6, 12, 18, 36\}$ steps, 1\,h to 6\,h) evaluates the sensitivity of the training loss objective when indexing offline power curve priors via lagged historical wind speed measurements in the absence of forward numerical weather predictions.

\subsection*{Quantification and Statistical Analysis}

\paragraph*{Operational Decision Risk Formulation}
To eliminate symmetric evaluation bias, performance was quantified using the non-negative asymmetric Pinball Loss function:
\begin{equation}
\mathcal{L}_q(y, \hat{y}) = \max\big( q(y - \hat{y}),\, (1-q)(\hat{y} - y) \big),
\end{equation}
where $q = \frac{1}{r+1} = \frac{1}{3}$ corresponds to an operational market penalty ratio $r=2$ (under-generation penalty twice as costly as spill). Total normalized static cost was defined as $\text{Cost}_{\text{static}} = 3 \times \mathcal{L}_{1/3}(y, \hat{y})$, with validation-set quantile post-hoc offset calibration $\delta^* = \arg\min_\delta \hat{\mathcal{L}}_{1/3}(y, \hat{y} + \delta)$.

\paragraph*{Circular Block Paired Bootstrap Inference}
Statistical significance of exploratory dynamic price model contrasts ($\Delta \text{Cost} = \text{Cost}_{\text{method}} - \text{Cost}_{\text{baseline}}$, Table~\ref{tab:paired_bootstrap}) was evaluated via circular block paired bootstrapping with 2,000 replications. Block length was set to $L = 1,008$ intervals (7 days) to rigorously preserve temporal autocorrelation and diurnal meteorological cycles. Two-sided 95\% percentile confidence intervals were constructed from the empirical bootstrap distribution, rejecting the null hypothesis of equal performance if and only if zero fell outside the confidence interval ($p < 0.05$). For the prospective 3D controlled ablations in Table~\ref{tab:controlled_ablation}, metric margins represent sample means and standard deviations across independent random initialization seeds.
